%------------------------------------------------------------------------------%

\usepackage{calculo_varias_variaveis-1}

\title{Máximos e Mínimos}

%------------------------------------------------------------------------------%
\begin{document}

\addtitle

%------------------------------------------------------------------------------%
\section{Extremos locais}

%------------------------------------------------------------------------------%
\begin{frame}{Extremos locais}
  Seja $f(x, y)$ definida em uma região $R$ que contém o ponto $(a, b)$

  \pause
  \begin{itemize}[<+->]
    \item \(f(a, b)\) é um valor \emph{máximo local} de $f$ se
          \[
            f(a, b) \geq f(x, y)
          \]
          para \((x, y)\) do domínio em um disco aberto centrado em \((a, b)\)
          \\[5mm]
    \item \(f(a, b)\) é um valor \emph{mínimo local} de $f$ se
          \[
            f(a, b) \leq f(x, y)
          \]
          para \((x, y)\) do domínio em um disco aberto centrado em \((a, b)\)
  \end{itemize}
\end{frame}

%------------------------------------------------------------------------------%
\begin{frame}{Teste da derivada de primeira ordem}
  Se $f(x, y)$ tiver um máximo ou mínimo local\\
  em um ponto interior \((a, b)\) do seu domínio

  \pause
  \vspace{0.1\baselineskip}
  e se as derivadas de primeira ordem existirem em \((a, b)\), então

  \pause
  \vspace{0.1\baselineskip}
  \[
    f_x(a, b) = 0
    \qquad
    f_y(a, b) = 0
  \]
\end{frame}

%------------------------------------------------------------------------------%
\begin{frame}{Ponto crítico}

  Um \emph{Ponto Crítico} de uma função \(f(x, y)\)

  \pause
  \vspace{0.1\baselineskip}
  é um ponto $(a, b)$, no interior do domínio de $f$, onde

  \pause
  \vspace{0.1\baselineskip}
  \emph{\(f_x(a, b) = 0\)} e \emph{\(f_y(a, b) =0\)  }
  \pause
  \tabto{70mm}
  \(\nabla f(a, b) = 0\)

  \pause
  \vspace{0.1\baselineskip}
  ou

  \emph{\(f_x(a, b)\) ou \(f_y(a, b)\) não exista}
  \pause
  \tabto{70mm}
  \(\nabla f(a, b)\) não existe

\end{frame}

%------------------------------------------------------------------------------%
\begin{frame}{Ponto de sela}
  Uma função diferenciável \(f(x, y)\) possui um \emph{ponto de sela} em \((a, b)\)

  \pause
  \vspace{0.1\baselineskip}
  se em todo disco aberto centrado em \((a, b)\)

  \pause
  \vspace{0.1\baselineskip}
  existirem pontos \((x, y)\) onde
  \[
    f(x, y) > f(a, b)
  \]
  \pause
  e existirem pontos \((x, y)\) onde
  \[
    f(x, y) < f(a, b)
  \]
\end{frame}

%------------------------------------------------------------------------------%
%------------------------------------------------------------------------------%
\addtocounter{exemplo}{1}
\begin{frame}{Exemplo \theexemplo}
  Encontre os pontos críticos e valores extremos locais da função
  \[
     f(x, y) = x^2 + y^2 - 4y + 9
  \]
\end{frame}

%------------------------------------------------------------------------------%
\begin{frame}{Exemplo \theexemplo}
  Derivadas parciais
  \pause
  \vspace{\baselineskip}
  \[
    \D{f}{x}
    \pause
    = \D{}{x}\left(x^2 + y^2 - 4y + 9\right)
    \pause
    = 2x
  \]
  \pause
  \vspace{\baselineskip}
  \[
    \D{f}{y}
    \pause
    = \D{}{y}\left(x^2 + y^2 - 4y + 9\right)
    \pause
    = 2y -4
  \]
\end{frame}

%------------------------------------------------------------------------------%
\begin{frame}{Exemplo \theexemplo}
  \onslide<+->{Pontos críticos}

  \vspace{0.1\baselineskip}
  \onslide<+->{As derivadas existem em todos os pontos do plano}

  \vspace{0.1\baselineskip}
  \onslide<+->{Plano tangente horizontal}
  \begin{columns}
  \column{0.33\linewidth}
    \begin{align*}
      \onslide<+->{f_x(x, y) & = 0 \\}
      \onslide<+->{2x        & = 0 \\}
      \onslide<+->{x         & = 0}
    \end{align*}
  \column{0.33\linewidth}
    \begin{align*}
      \onslide<+->{f_y(x, y) & = 0 \\}
      \onslide<+->{2y -4     & = 0 \\}
      \onslide<+->{y         & = 2}
    \end{align*}
    \column{0.33\linewidth}
    \onslide<+->{Ponto critico é \((0, 2)\)}
  \end{columns}
\end{frame}

%------------------------------------------------------------------------------%
\begin{frame}{Exemplo \theexemplo}
  Onde a função vale
  \[
    f(0, 2)
    \pause
    = \left(x^2 + y^2 - 4y + 9\right)\evalat{(0, 2)}{}
    \pause
    = 0 + 4 - 8 + 9
    \pause
    = 5
  \]

  \pause
  \vspace{\baselineskip}
  Analisando a função percebemos que 5 é o menor valor que ela assume,
  portanto é um ponto de mínimo local (e nesse caso global também)
\end{frame}

%------------------------------------------------------------------------------%


%------------------------------------------------------------------------------%
\section{Teste da derivada segunda}

%------------------------------------------------------------------------------%
\begin{frame}{Hessiana}
  \emph{Matriz Hessiana}
  \[
    \arraycolsep=2mm
    \def\arraystretch{1.2}
    H(x, y) = \begin{bmatrix}
      \null\;f_{xx}(x, y)\;\null & \null\;f_{xy}(x, y)\;\null \\
      f_{yx}(x, y) & f_{yy}(x, y)
    \end{bmatrix}
  \]

  \pause
  \vspace{\baselineskip}
  \emph{Discriminante (ou hessiano)}
  \[
    D(x, y) = \det H(x, y) \;=\; f_{xx}(x, y)\,f_{yy}(x, y) \;-\; f_{xy}^2(x, y)
  \]
\end{frame}

%------------------------------------------------------------------------------%
\begin{frame}{Teste da derivada segunda}
  Seja $f(x, y)$ diferenciável
  \pause
  e $(a, b)$ um ponto crítico onde
  \emph{\(
    f_x(a, b) = f_y(a, b) = 0
  \)}

  \pause
  \begin{enumerate}[<+->]
    \setlength{\itemsep}{4mm}
    \item se \emph{\(\det H(a, b) = 0\)} \tabto{40mm}
      o teste é inconclusivo
    \item se \emph{\(\det H(a, b) < 0\)} \tabto{40mm}
      $f$ tem um \emph{ponto de sela} em $(a, b)$
    \item se \emph{\(\det H(a, b) > 0\)} \tabto{40mm}
      testamos a derivada segunda em $x$ (ou $y$) \\[3mm]
      \begin{enumerate}
        \setlength{\itemsep}{3mm}
        \item \normalsize\; se \emph{\(f_{xx}(a, b) < 0\)} \tabto{35mm}
          $f$ tem um \emph{máximo local} em $(a, b)$
        \item \normalsize\; se \emph{\(f_{xx}(a, b) > 0\)} \tabto{35mm}
          $f$ tem um \emph{mínimo local} em $(a, b)$
      \end{enumerate}
  \end{enumerate}
\end{frame}

%------------------------------------------------------------------------------%
\section{Exemplos}

%------------------------------------------------------------------------------%
\begin{question}
  Encontre e classifique os pontos críticos da função
  \[
    f(x, y) = xy -x^2 - y^2 -2x -2y +4
  \]
  Avalie os valores extremos locais de $f$
\end{question}

%------------------------------------------------------------------------------%
\begin{resolution}{Derivadas parciais}
  \[
    \D{f}{x}
    \pause
    = \D{}{x}\left(xy -x^2 - y^2 -2x -2y +4\right)
    \pause
    = y -2x -2
  \]
  \pause
  \vspace{\baselineskip}
  \[
    \D{f}{y}
    \pause
    = \D{}{y}\left(xy -x^2 - y^2 -2x -2y +4\right)
    \pause
    = x -2y -2
  \]
\end{resolution}

%------------------------------------------------------------------------------%
\begin{resolution}{Pontos críticos}
  As derivadas existem em todos os pontos do plano

  \pause
  \vspace{\baselineskip}
  Buscando os pontos onde plano tangente é horizontal
\end{resolution}

%------------------------------------------------------------------------------%
\begin{resolution}{Resolvendo o sistema}
\begin{columns}
\column{0.5\linewidth}
  \begin{align*}
    \onslide<+->{f_x(x, y) & = y -2x -2 = 0 \\}
    \onslide<+->{f_y(x, y) & = x -2y -2 = 0}
  \end{align*}

  \vspace{0.1\baselineskip}
  \begin{align*}
    \onslide<+->{x -2y -2 & = 0      \\}
    \onslide<+->{x        & = 2y + 2}
  \end{align*}

\column{0.5\linewidth}
  \begin{align*}
    \onslide<+->{y -2x -2         & = 0  \\}
    \onslide<+->{y -2(2y + 2) - 2 & = 0  \\}
    \onslide<+->{y -4y + -6       & = 0  \\}
    \onslide<+->{-3y              & = 6  \\}
    \onslide<+->{y                & = -2}
  \end{align*}

  \[
    \onslide<+->{x = 2y + 2}
    \onslide<+->{= 2(-2) + 2}
    \onslide<+->{= -2}
  \]

  \onslide<+->{Ponto crítico $(-2, -2)$}
\end{columns}
\end{resolution}

%------------------------------------------------------------------------------%
\begin{resolution}{Calculando as derivadas segundas}
  \[
    \pause
    \DS{f}{x}
    \pause
    = \D{}{x}\left(y -2x -2\right)
    \pause
    = -2
  \]
  \[
    \pause
    \DS{f}{y}
    \pause
    = \D{}{y}\left(x -2y -2\right)
    \pause
    = -2
  \]
  \[
    \pause
    \DM{f}{x}{y}
    \pause
    = \D{}{x}\left(x -2y -2\right)
    \pause
    = 1
  \]

  \pause
  Discriminante
  \[
    \pause
    D(x, y)
    \pause
    = f_{xx}(x, y) f_{yy}(x, y) - f^2_{xy}(x, y)
    \pause
    = (-2)(-2) - 1
    \pause
    = 3
  \]
  \[
  \pause
    D(-2, -2)
    \pause
    = 3
  \]
\end{resolution}

%------------------------------------------------------------------------------%
\begin{resolution}{Caracterização do ponto}
  Como \(D(-2, -2)=3>0\)
  \pause
  e  \(f_{xx}(-2, -2) = -2 < 0\)

  \pause
  \vspace{\baselineskip}
  O ponto \((-2, -2)\) é um ponto de máximo local

  \pause
  \vspace{\baselineskip}
  O valor da função neste máximo local é
  \[
    f(-2, -2)
    \pause
    = \left(xy -x^2 - y^2 -2x -2y +4\right)\evalat{(-2, -2)}{}
  \]
  \[
    \hphantom{f(-2, -2)}
    \pause
    = (-2)(-2) - (-2)^2 - (-2)^2 -2(-2) -2(-2) + 4
  \]
  \[
    \hphantom{f(-2, -2)}
    \pause
    = 4 -4 -4 +4 +4 + 4
  \]
  \[
    \hphantom{f(-2, -2)}
    \pause
    = 8
  \]
\end{resolution}

%------------------------------------------------------------------------------%

%------------------------------------------------------------------------------%
\addtocounter{exemplo}{1}
\begin{frame}{Exemplo \theexemplo}
  Encontre e classifique os pontos críticos da função
  \[
    f(x, y) = 3y^2 -2y^3 -3x^2 +6xy
  \]
  Avalie os valores extremos locais de $f$
\end{frame}

%------------------------------------------------------------------------------%
\begin{frame}{Exemplo \theexemplo}
  Derivadas parciais
  \pause
  \vspace{\baselineskip}
  \[
    \D{f}{x}
    \pause
    = \D{}{x}\left(3y^2 -2y^3 -3x^2 +6xy\right)
    \pause
    = -6x +6y
    \pause
    = 6y -6x
  \]
  \pause
  \vspace{\baselineskip}
  \[
    \D{f}{y}
    \pause
    = \D{}{y}\left(3y^2 -2y^3 -3x^2 +6xy\right)
    \pause
    = 6y -6y^2 +6x
  \]
\end{frame}

%------------------------------------------------------------------------------%
\begin{frame}{Exemplo \theexemplo}
  Pontos críticos

  \pause
  \vspace{\baselineskip}
  As derivadas existem em todos os pontos do plano
\end{frame}

%------------------------------------------------------------------------------%
\begin{frame}{Exemplo \theexemplo}
\begin{columns}
\column{0.5\linewidth}
  \begin{align*}
    \onslide<+->{f_x(x, y) & = 6y -6x = 0 \\}
    \onslide<+->{f_y(x, y) & = 6y -6y^2 +6x = 0}
  \end{align*}
  \begin{align*}
    \onslide<+->{6y -6x & = 0 \\}
    \onslide<+->{x      & = y}
  \end{align*}
  \column{0.5\linewidth}
  \begin{align*}
    \onslide<+->{6y -6y^2 +6x & = 0 \\}
    \onslide<+->{6y^2 -6y -6x & = 0 \\}
    \onslide<+->{6y^2 -6y -6y & = 0 \\}
    \onslide<+->{6y^2 - 12y   & = 0 \\}
    \onslide<+->{6y(y-2)      & = 0}
  \end{align*}

  \vspace{\baselineskip}
  \onslide<+->{Portanto \(y=0\) ou \(y=2\)}

  \vspace{\baselineskip}
  \onslide<+->{Os pontos críticos são \((0, 0)\) e \((2, 2)\)}
\end{columns}
\end{frame}

%------------------------------------------------------------------------------%
\begin{frame}{Exemplo \theexemplo}
\begin{columns}[t]
\column{0.5\linewidth}
  \onslide<+->{Calculando as derivadas segundas}
  \[
    \onslide<+->{\DS{f}{x}}
    \onslide<+->{= \D{}{x}\left(6y -6x\right)}
    \onslide<+->{= -6}
  \]

  ~
  \[
    \onslide<+->{\DS{f}{y}}
    \onslide<+->{= \D{}{y}\left(6y -6y^2 +6x\right)}
    \onslide<+->{= 6 -12y}
  \]

  ~
  \[
    \onslide<+->{\DM{f}{x}{y}}
    \onslide<+->{= \DM{f}{y}{x}}
    \onslide<+->{= \D{}{y}\left(6y -6x\right)}
    \onslide<+->{= 6}
  \]
\column{0.5\linewidth}
  \onslide<+->{Discriminante}
  \begin{align*}
    \onslide<+->{D(x, y)}
    \onslide<+->{& = f_{xx}(x, y) f_{yy}(x, y) - f^2_{xy}(x, y) \\}
    \onslide<+->{& = (-6)(6-12y) - 6^2 \\}
    \onslide<+->{& = -36 +72y - 36 \\}
    \onslide<+->{& = 72(y-1)}
  \end{align*}
\end{columns}
\end{frame}

%------------------------------------------------------------------------------%
\begin{frame}{Exemplo \theexemplo}
  Analisando o ponto \((0, 0)\)

  \pause
  \vspace{\baselineskip}
  Como
  \[
    \pause
    D(0, 0)
    \pause
    = 72(0-1)
    \pause
    = -72
    \pause
    \;<\; 0
  \]
  \pause
  esse ponto é um ponto de sela
\end{frame}

%------------------------------------------------------------------------------%
\begin{frame}{Exemplo \theexemplo}
  \onslide<+->{Analisando o ponto \((2, 2)\)}

  \onslide<+->{Como}
  \[
    \onslide<+->{D(2, 2)}
    \onslide<+->{= 72(2-1)}
    \onslide<+->{= 72}
    \onslide<+->{> 0}
    \qquad
    \onslide<+->{f_{xx}}
    \onslide<+->{= -6}
    \onslide<+->{< 0}
  \]
  \onslide<+->{o ponto \((2, 2)\) é um máximo local com valor}
  \begin{align*}
    \onslide<+->{f(2, 2)}
    \onslide<+->{& = \left(3y^2 -2y^3 -3x^2 +6xy\right)\evalat{(2, 2)}{} \\}
    \onslide<+->{& = 3\times 2^2 -2 \times 2^3 -3\times 2^2 + 6\times 2 \times 2\\}
    \onslide<+->{& = 12 -16 -12 + 24 \\}
    \onslide<+->{& = 8}
  \end{align*}
\end{frame}

%------------------------------------------------------------------------------%

%------------------------------------------------------------------------------%
\addtocounter{exemplo}{1}
\begin{frame}{Exemplo \theexemplo}
  Seja
  \[
    f(x, y) = e^x \cos(y)
  \]
  encontre os pontos críticos de $f$ e classifique-os
\end{frame}

%------------------------------------------------------------------------------%
\begin{frame}{Exemplo \theexemplo}
  \onslide<+->{Derivadas primeiras}
  \begin{align*}
    \onslide<+->{f_x(x, y) & = e^x \cos(y) &}
    \onslide<+->{f_y(x, y) & = - e^x \sin(y)}
  \end{align*}

  \onslide<+->{Pontos onde as derivadas são nulas}
  \begin{columns}
  \column{0.5\linewidth}
    \begin{align*}
      \onslide<+->{f_x(x, y)  & = 0                                 \\}
      \onslide<+->{e^x\cos(y) & = 0                                 \\}
      \onslide<+->{\cos(y)    & = 0                                 \\}
      \onslide<+->{y          & = \frac{\pi}{2} + n\pi \quad n \in \Z }
    \end{align*}
  \column{0.5\linewidth}
    \begin{align*}
      \onslide<+->{f_y(x, y)   & = 0                 \\}
      \onslide<+->{-e^x\sin(y) & = 0                 \\}
      \onslide<+->{\sin(y)     & = 0                 \\}
      \onslide<+->{y           & = n\pi \quad n \in \Z }
    \end{align*}
  \end{columns}
\end{frame}

%------------------------------------------------------------------------------%
\begin{frame}{Exemplo \theexemplo}
  Não existe solução simultânea

  \pause
  e as derivadas existem em todos os pontos de $\R^2$,

  \pause
  portanto, não há pontos críticos
\end{frame}

%------------------------------------------------------------------------------%

%------------------------------------------------------------------------------%
\begin{question}
  Seja
  \[
    f(x, y) = \ln(x^2 + y^2 + 1)
  \]
  encontre os pontos críticos de $f$ e classifique-os
\end{question}

%------------------------------------------------------------------------------%
\begin{resolution}{Derivadas primeiras}
  \begin{align*}
    \onslide<+->{f_x(x, y) & }
    \onslide<+->{= \frac{1}{x^2 + y^2 + 1} \,\D{}{x}(x^2 + y^2 + 1)}
    \onslide<+->{= \frac{2x}{x^2 + y^2 + 1} \\}
    \onslide<+->{f_y(x, y) & }
    \onslide<+->{= \frac{1}{x^2 + y^2 + 1} \,\D{}{y}(x^2 + y^2 + 1)}
    \onslide<+->{= \frac{2y}{x^2 + y^2 + 1}}
  \end{align*}
  \onslide<+->{Pontos onde as derivadas são nulas}
  \begin{columns}
    \column{0.33\linewidth}
    \begin{align*}
      \onslide<+->{f_x(x, y)                & = 0 \\}
      \onslide<+->{\frac{2x}{x^2 + y^2 + 1} & = 0 \\}
      \onslide<+->{x                        & = 0 }
    \end{align*}
    \column{0.33\linewidth}
    \begin{align*}
      \onslide<+->{f_y(x, y)                & = 0 \\}
      \onslide<+->{\frac{2y}{x^2 + y^2 + 1} & = 0 \\}
      \onslide<+->{y                        & = 0 }
    \end{align*}
    \column{0.33\linewidth}
    \onslide<+->{Ponto crítico $(0, 0)$}
  \end{columns}
\end{resolution}

%------------------------------------------------------------------------------%
\begin{resolution}{Derivadas segundas}
  \begin{align*}
    \onslide<+->{f_{xx}(x, y)}
    \onslide<+->{& = \D{}{x}\left(\frac{2x}{x^2 + y^2 + 1}\right) \\}
    \onslide<+->{& = \frac{\displaystyle\D{}{x}(2x)(x^2 + y^2 + 1) - 2x\D{}{x}(x^2 + y^2 + 1)}{(x^2 + y^2 + 1)^2} \\}
    \onslide<+->{& = \frac{2(x^2 + y^2 + 1) - 2x(2x)}{(x^2 + y^2 + 1)^2} \\}
    \onslide<+->{& = \frac{2(y^2 + 1 - x^2)}{(x^2 + y^2 + 1)^2} }
  \end{align*}
\end{resolution}

%------------------------------------------------------------------------------%
\begin{resolution}{Derivadas segundas}
  \begin{align*}
    \onslide<+->{f_{yy}(x, y)}
    \onslide<+->{& = \D{}{y}\left(\frac{2y}{x^2 + y^2 + 1}\right) \\}
    \onslide<+->{& = \frac{\displaystyle\D{}{y}(2y)(x^2 + y^2 + 1) - 2y\D{}{y}(x^2 + y^2 + 1)}{(x^2 + y^2 + 1)^2} \\}
    \onslide<+->{& = \frac{2(x^2 + y^2 + 1) - 2y(2y)}{(x^2 + y^2 + 1)^2} \\}
    \onslide<+->{& =  \frac{2(x^2 + 1 - y^2)}{(x^2 + y^2 + 1)^2}}
  \end{align*}
\end{resolution}

%------------------------------------------------------------------------------%
\begin{resolution}{Derivadas segundas}
  \begin{align*}
    \onslide<+->{f_{xy}(x, y)}
    \onslide<+->{& = \D{}{y}\left(\frac{2x}{x^2 + y^2 + 1}\right) \\}
    \onslide<+->{& = \frac{\displaystyle\D{}{y}(2x)(x^2 + y^2 + 1) - 2x\D{}{y}(x^2 + y^2 + 1)}{(x^2 + y^2 + 1)^2} \\}
    \onslide<+->{& = \frac{0 - 2x(2y)}{(x^2 + y^2 + 1)^2} \\}
    \onslide<+->{& =  \frac{-4xy}{(x^2 + y^2 + 1)^2}}
  \end{align*}
\end{resolution}

%------------------------------------------------------------------------------%
\begin{resolution}{Derivadas segundas}
  \begin{align*}
    \onslide<+->{f_{xx}(x, y) & = \frac{2(y^2 + 1 - x^2)}{(x^2 + y^2 + 1)^2} &}
    \onslide<+->{f_{xx}(0, 0) & = 2 \\}
    \onslide<+->{f_{yy}(x, y) & = \frac{2(x^2 + 1 - y^2)}{(x^2 + y^2 + 1)^2} &}
    \onslide<+->{f_{yy}(0, 0) & = 2 \\}
    \onslide<+->{f_{xy}(x, y) & = \frac{-4xy}{(x^2 + y^2 + 1)^2}             &}
    \onslide<+->{f_{xy}(0, 0) & = 0}
  \end{align*}
\end{resolution}

%------------------------------------------------------------------------------%
\begin{resolution}{Discriminante}
  \begin{align*}
    \onslide<+->{D(0, 0)}
    \onslide<+->{& = f_{xx}(0, 0) f_{yy}(0, 0) - f^2_{xy}(0, 0) \\}
    \onslide<+->{& = 2 \times 2 - 0^2 \\}
    \onslide<+->{& = 4 }
    \onslide<+->{\;>\; 0}
  \end{align*}

  \[
    \onslide<+->{f_{xx}(0, 0)}
    \onslide<+->{= 2}
    \onslide<+->{> 0}
  \]

  \onslide<+->{O ponto $(0, 0)$ é um mínimo local}
\end{resolution}

%------------------------------------------------------------------------------%

%------------------------------------------------------------------------------%
\begin{question}
  Seja
  \[
    f(x, y) = x e^{-x^2 - y^2}
  \]
  Encontre os pontos críticos de $f$ e classifique-os
\end{question}

%------------------------------------------------------------------------------%
\begin{resolution}{Derivadas primeiras}
  \begin{align*}
    \onslide<+->{f_x(x, y)}
    \onslide<+->{& = \D{}{x}\left(x e^{-x^2 - y^2}\right)                             \\}
    \onslide<+->{& = \D{x}{x}e^{-x^2 - y^2} + x \D{}{x}e^{-x^2 - y^2}                 \\}
    \onslide<+->{& = e^{-x^2 - y^2} + x e^{-x^2 - y^2} \D{}{x}\left(-x^2 - y^2\right) \\}
    \onslide<+->{& = e^{-x^2 - y^2} + x e^{-x^2 - y^2} \left(-2x\right)               \\}
    \onslide<+->{& = \left(1 - 2x^2\right)e^{-x^2 - y^2} }
  \end{align*}
\end{resolution}

%------------------------------------------------------------------------------%
\begin{resolution}{Derivadas primeiras}
  \begin{align*}
    \onslide<+->{f_y(x, y)}
    \onslide<+->{& = \D{}{y}\left(x e^{-x^2 - y^2}\right)            \\}
    \onslide<+->{& = x \D{}{y}e^{-x^2 - y^2}                         \\}
    \onslide<+->{& = x e^{-x^2 - y^2} \D{}{y}\left(-x^2 - y^2\right) \\}
    \onslide<+->{& = x e^{-x^2 - y^2} \left(-2y\right)               \\}
    \onslide<+->{& = -2xy e^{-x^2 - y^2}}
  \end{align*}
\end{resolution}

%------------------------------------------------------------------------------%
\begin{resolution}{Pontos críticos}
  As derivadas existem em todos os pontos do plano

  \pause
  Pontos com plano tangente horizontal
\end{resolution}

%------------------------------------------------------------------------------%
\begin{resolution}{Sistema}
\begin{columns}[t]
\column{0.5\linewidth}
  \[
    \begin{cases}
      f_y(x, y) = 0 \\[3mm]
      f_x(x, y) = 0
    \end{cases}
  \]

  \pause
  \vspace{\baselineskip}
  \[
    \begin{cases}
      \left(1 - 2x^2\right)e^{-x^2 - y^2} = 0 \\[3mm]
      -2xy e^{-x^2 - y^2} = 0
    \end{cases}
  \]
\column{0.5\linewidth}
  \pause
    \[
    \begin{cases}
      1 - 2x^2 = 0 \\[3mm]
      xy = 0
    \end{cases}
  \]

  \pause
  \vspace{\baselineskip}
  \[
    xy = 0
  \]
  \[
    \pause
    x = 0
    \pause
    \;\text{ ou }\; y = 0
  \]
\end{columns}
\end{resolution}

%------------------------------------------------------------------------------%
\begin{resolution}{Sistema}
\begin{columns}[t]
\column{0.5\linewidth}
  \onslide<+->{Se $x=0$}
  \begin{align*}
    \onslide<+->{f_x(0, y)       & = 0 \\}
    \onslide<+->{1 - 2\times 0^2 & = 0 \\}
    \onslide<+->{1               & = 0}
  \end{align*}
  \onslide<+->{Não existe solução}
\column{0.5\linewidth}
  \onslide<+->{Se $y=0$}
  \begin{align*}
    \onslide<+->{f_x(x, 0) & = 0                     \\}
    \onslide<+->{1 - 2x^2  & = 0                     \\}
    \onslide<+->{2x^2      & = 1                     \\}
    \onslide<+->{x         & = \frac{\pm1}{\sqrt{2}}}
  \end{align*}
\end{columns}

  \onslide<+->{Pontos críticos
  \(\left( \frac{-1}{\sqrt{2}}, 0 \right)\) e
  \(\left( \frac{ 1}{\sqrt{2}}, 0 \right)\)}
\end{resolution}

%------------------------------------------------------------------------------%
\begin{resolution}{Derivadas segundas}
  \begin{align*}
    \onslide<+->{f_{xx}(x, y)}
    \onslide<+->{& = \D{}{x}\left(\left(1 - 2x^2\right)e^{-x^2 - y^2}\right) \\}
    \onslide<+->{& = \D{}{x}\left(1 - 2x^2\right)e^{-x^2 - y^2} + \left(1 - 2x^2\right)\D{}{x}e^{-x^2 - y^2} \\}
    \onslide<+->{& = -4x e^{-x^2 - y^2} + \left(1 - 2x^2\right)e^{-x^2 - y^2}\D{}{x}\left(-x^2 - y^2\right) \\}
    \onslide<+->{& = -4x e^{-x^2 - y^2} + \left(1 - 2x^2\right)e^{-x^2 - y^2}\left(-2x\right) \\}
    \onslide<+->{& = \left[-4x + \left(1 - 2x^2\right)\left(-2x\right)\right] e^{-x^2 - y^2} \\}
    \onslide<+->{& = 2x\left(2x^2 - 3\right) e^{-x^2 - y^2}}
  \end{align*}
\end{resolution}

%------------------------------------------------------------------------------%
\begin{resolution}{Derivadas segundas}
  \begin{align*}
    \onslide<+->{f_{yy}(x, y)}
    \onslide<+->{& = \D{}{y}\left(-2xy e^{-x^2 - y^2}\right) \\}
    \onslide<+->{& = -2x \left(\D{y}{y}e^{-x^2 - y^2} + y\D{}{y}e^{-x^2 - y^2}\right) \\}
    \onslide<+->{& = -2x \left(e^{-x^2 - y^2} + ye^{-x^2 - y^2}\D{}{y}\left(-x^2 - y^2\right)\right) \\}
    \onslide<+->{& = -2x \left(e^{-x^2 - y^2} + ye^{-x^2 - y^2}(-2y)\right) \\}
    \onslide<+->{& = -2x \left(1 -2y^2\right)e^{-x^2 - y^2} \\}
    \onslide<+->{& =  2x \left(2y^2 - 1\right) e^{-x^2 - y^2}}
  \end{align*}
\end{resolution}

%------------------------------------------------------------------------------%
\begin{resolution}{Derivadas segundas}
  \begin{align*}
    \onslide<+->{f_{xy}(x, y)}
    \onslide<+->{& = \D{}{y}\left(\left(1 - 2x^2\right)e^{-x^2 - y^2}\right) \\}
    \onslide<+->{& = \left(1 - 2x^2\right)\D{}{y}e^{-x^2 - y^2} \\}
    \onslide<+->{& = \left(1 - 2x^2\right)e^{-x^2 - y^2}\D{}{y}\left(-x^2 - y^2\right) \\}
    \onslide<+->{& = \left(1 - 2x^2\right)e^{-x^2 - y^2}\left(-2y\right) \\}
    \onslide<+->{& = 2y\left(2x^2-1\right)e^{-x^2 - y^2}}
  \end{align*}
\end{resolution}

%------------------------------------------------------------------------------%
\begin{resolution}{Derivadas segundas}
  Derivadas segundas
  \begin{align*}
    f_{xx}(x, y) & = 2x \left(2x^2 - 3\right) e^{-x^2 - y^2} \\
    f_{yy}(x, y) & = 2x \left(2y^2 - 1\right) e^{-x^2 - y^2} \\
    f_{xy}(x, y) & = 2y \left(2x^2 - 1\right) e^{-x^2 - y^2}
  \end{align*}
\end{resolution}

%------------------------------------------------------------------------------%
\begin{resolution}{Derivadas segundas}
  \begin{align*}
    \onslide<+->{f_{xx}\left(\frac{1}{\sqrt{2}}, 0\right)}
    \onslide<+->{& = 2\left(2x^3 - 3x\right) e^{-x^2 - y^2} \evalat{\left(\sfrac{1}{\sqrt{2}}, 0\right)}{}  \\}
    \onslide<+->{& = 2\left(2\left(\frac{1}{\sqrt{2}}\right)^3 - \frac{3}{\sqrt{2}}\right)
                     e^{-\left(\sfrac{1}{\sqrt{2}}\right)^2}  \\}
    \onslide<+->{& = 2\left(\frac{2}{2\sqrt{2}} - \frac{3}{\sqrt{2}}\right) e^{-\sfrac{1}{2}}  \\}
    \onslide<+->{& = 2\frac{1-3}{\sqrt{2}} e^{-\sfrac{1}{2}} \\}
    \onslide<+->{& = -2\sqrt{2} e^{-\sfrac{1}{2}}}
  \end{align*}
\end{resolution}

%------------------------------------------------------------------------------%
\begin{resolution}{Derivadas segundas}
  \begin{align*}
    \onslide<+->{f_{yy}\left(\frac{1}{\sqrt{2}}, 0\right)}
    \onslide<+->{& = 2x \left(2y^2 - 1\right) e^{-x^2 - y^2} \evalat{\left(\sfrac{1}{\sqrt{2}}, 0\right)}{} \\}
    \onslide<+->{& = \frac{2}{\sqrt{2}} (-1) e^{-\sfrac{1}{2}}  \\}
    \onslide<+->{& = -\sqrt{2} e^{-\sfrac{1}{2}}}
  \end{align*}
\end{resolution}

%------------------------------------------------------------------------------%
\begin{resolution}{Derivadas segundas}
  \[
    f_{xy}\left(\frac{1}{\sqrt{2}}, 0\right)
    \pause
    = 2y\left(2x^2-1\right)e^{-x^2 - y^2} \evalat{\left(\sfrac{1}{\sqrt{2}}, 0\right)}{}
    \pause
    = 0
  \]
\end{resolution}

%------------------------------------------------------------------------------%
\begin{resolution}{Discriminante}
  \begin{align*}
    \onslide<+->{D\left(\frac{1}{\sqrt{2}}, 0\right)}
    \onslide<+->{= \det H\left(\frac{1}{\sqrt{2}}, 0\right)}
    \onslide<+->{& = f_{xx}\left(\frac{1}{\sqrt{2}}, 0\right)f_{yy}\left(\frac{1}{\sqrt{2}}, 0\right)
      - f^2_{xy}\left(\frac{1}{\sqrt{2}}, 0\right) \\}
    \onslide<+->{& = \left(-2\sqrt{2} e^{-\sfrac{1}{2}}\right)\left(-\sqrt{2} e^{-\sfrac{1}{2}}\right) - 0 \\}
    \onslide<+->{& = 4 e^{-1}}
    \onslide<+->{= \frac{4}{e}}
    \onslide<+->{> 0}
  \end{align*}

  \[
    \onslide<+->{f_{xx}\left(\frac{1}{\sqrt{2}}, 0\right)}
    \onslide<+->{= -2\sqrt{2} e^{-\sfrac{1}{2}}}
    \onslide<+->{< 0}
  \]

  \onslide<+->{O ponto \(\left(\frac{1}{\sqrt{2}}, 0\right)\) é um máximo local}
\end{resolution}

%------------------------------------------------------------------------------%
\begin{resolution}{Derivadas segundas}
  \begin{align*}
    \onslide<+->{f_{xx}\left(\frac{-1}{\sqrt{2}}, 0\right)}
    \onslide<+->{& = 2\left(2x^3 - 3x\right) e^{-x^2 - y^2} \evalat{\left(-\sfrac{1}{\sqrt{2}}, 0\right)}{}  \\}
    \onslide<+->{& = 2\left(2\left(\frac{-1}{\sqrt{2}}\right)^3 + \frac{3}{\sqrt{2}}\right)
                     e^{-\left(-\sfrac{1}{\sqrt{2}}\right)^2}  \\}
    \onslide<+->{& = 2\left(\frac{-2}{2\sqrt{2}} + \frac{3}{\sqrt{2}}\right) e^{-\sfrac{1}{2}}  \\}
    \onslide<+->{& = 2\frac{-1+3}{\sqrt{2}} e^{-\sfrac{1}{2}} \\}
    \onslide<+->{& = 2\sqrt{2} e^{-\sfrac{1}{2}}}
  \end{align*}
\end{resolution}

%------------------------------------------------------------------------------%
\begin{resolution}{Derivadas segundas}
  \begin{align*}
    \onslide<+->{f_{yy}\left(\frac{-1}{\sqrt{2}}, 0\right)}
    \onslide<+->{& = 2x \left(2y^2 - 1\right) e^{-x^2 - y^2} \evalat{\left(-\sfrac{1}{\sqrt{2}}, 0\right)}{} \\}
    \onslide<+->{& = \frac{-2}{\sqrt{2}} (-1) e^{-\sfrac{1}{2}}  \\}
    \onslide<+->{& = \sqrt{2} e^{-\sfrac{1}{2}}}
  \end{align*}
\end{resolution}

%------------------------------------------------------------------------------%
\begin{resolution}{Derivadas segundas}
  \[
    f_{xy}\left(\frac{-1}{\sqrt{2}}, 0\right)
    \pause
    = 2y\left(2x^2-1\right)e^{-x^2 - y^2} \evalat{\left(-\sfrac{1}{\sqrt{2}}, 0\right)}{}
    \pause
    = 0
  \]
\end{resolution}

%------------------------------------------------------------------------------%
\begin{resolution}{Discriminante}
  \begin{align*}
    \onslide<+->{D\left(\frac{-1}{\sqrt{2}}, 0\right)}
    \onslide<+->{= \det H\left(\frac{-1}{\sqrt{2}}, 0\right)}
    \onslide<+->{& = f_{xx}  \left(\frac{-1}{\sqrt{2}}, 0\right)
                     f_{yy}  \left(\frac{-1}{\sqrt{2}}, 0\right)
                   - f^2_{xy}\left(\frac{-1}{\sqrt{2}}, 0\right) \\}
    \onslide<+->{& = \left(2\sqrt{2} e^{-\sfrac{1}{2}}\right)\left(\sqrt{2} e^{-\sfrac{1}{2}}\right) - 0 \\}
    \onslide<+->{& = 4 e^{-1}}
    \onslide<+->{= \frac{4}{e}}
    \onslide<+->{> 0}
  \end{align*}

  \[
    \onslide<+->{f_{xx}\left(\frac{-1}{\sqrt{2}}, 0\right)}
    \onslide<+->{= 2\sqrt{2} e^{-\sfrac{1}{2}}}
    \onslide<+->{> 0}
  \]

  \onslide<+->{O ponto \(\left(\frac{-1}{\sqrt{2}}, 0\right)\) é um mínimo local}
\end{resolution}

%------------------------------------------------------------------------------%

%------------------------------------------------------------------------------%
\begin{question}
  Considerando a função
  \(
    f(x, y) = x^3 + 3xy + y^3
  \)

  a) Calcule o gradiente de $f$

  b) Calcule a hessiana de $f$

  c) Encontre todos os pontos críticos de $f$

  d) Classifique cada ponto crítico de $f$
\end{question}

%------------------------------------------------------------------------------%
\begin{resolution}{a}
  Precisamos das derivadas parciais de $f$
  \[
    \D{f}{x}
    = \D{}{x}\left[x^3 + 3xy + y^3\right]
    = 3x^2 + 3y
  \]
  \[
    \D{f}{y}
    = \D{}{y}\left[x^3 + 3xy + y^3\right]
    = 3x + 3y^2
  \]
  então
  \[
    \nabla f = \vetor{3x^2 + 3y}{3x + 3y^2}
  \]
\end{resolution}

%------------------------------------------------------------------------------%
\begin{resolution}{b}
  Precisamos das derivadas parciais de segunda ordem de $f$
  \[
    \DS{f}{x}
    = \D{}{x}\left[3x^2 + 3y\right]
    = 6x
  \]
  \[
    \D{f}{y}
    = \D{}{y}\left[3x + 3y^2\right]
    = 6y
  \]
  \[
    \DM{f}{x}{y}
    = \D{}{x}\D{f}{y}
    = \D{}{x}\left[3x + 3y^2\right]
    = 3
  \]
  então
  \[
    H = \left(\begin{array}{cc}
      6x & 3 \\
      3 & 6y
    \end{array}\right)
  \]
\end{resolution}

%------------------------------------------------------------------------------%
\begin{resolution}{c}
  A função é um polinômio, então possui derivadas em todos os pontos do plano.
  Assim os pontos críticos são apenas os pontos onde as derivadas parciais são zero,
  \(\nabla f = 0\)
  \[
    3x^2 + 3y = 0
    \qquad\text{e}\qquad
    3x + 3y^2 = 0
  \]
  ou, simplificando,
  \[
    x^2 + y = 0
    \qquad\text{e}\qquad
    x + y^2 = 0
  \]
\end{resolution}

%------------------------------------------------------------------------------%
\begin{resolution}{c}
  Isolando $y$ na primeira equação, \(y = -x^2\), e substituindo na segunda, temos
  \begin{align*}
    x + y^2                 & = 0 \\
    x + \left(-x^2\right)^2 & = 0 \\
    x + x^4                 & = 0 \\
    x(1+ x^3)               & = 0
  \end{align*}
  As soluções dessa equação são $x=0$ ou $x-1$.
  Se $x=0$ temos $y=0$ e
  se $x=-1$ temos $y=-1$.
  Portanto, os pontos críticos são
  \[
    (x_1, y_1) = (0, 0)
    \qquad\text{e}\qquad
    (x_2, y_2) = (-1, -1)
  \]
\end{resolution}

%------------------------------------------------------------------------------%
\begin{resolution}{d}
  Para classificar os pontos críticos precisamos avaliar o discriminante nos
  pontos críticos.
\end{resolution}

%------------------------------------------------------------------------------%
\begin{resolution}{d}
  Considerando o ponto \((x_1, y_1) = (0, 0)\)
  \[
    f_{xx}(0, 0) = 0
  \]
  \[
    f_{yy}(0, 0) = 0
  \]
  \[
    f_{xy}(0, 0) = 3
  \]
  \[
    D_1
    = f_{xx}(0, 0)f_{yy}(0, 0) - f^2_{xy}(0, 0)
    = 0 \times 0 - 3^2
    = -9 < 0
  \]
  Portanto, o ponto $(0, 0)$ é um ponto de sela.
\end{resolution}

%------------------------------------------------------------------------------%
\begin{resolution}{d}
  Considerando o ponto \((x_2, y_2) = (-1, -1)\)
  \[
    f_{xx}(-1, -1) = -6
  \]
  \[
    f_{yy}(-1, -1) = -6
  \]
  \[
    f_{xy}(-1, -1) = 3
  \]
  \[
    D_2
    = f_{xx}(-1, -1)f_{yy}(-1, -1) - f^2_{xy}(-1, -1)
    = (-6)(-6) - 3^2
    = 36 - 9
    = 25 > 0
  \]
  Portanto, o ponto $(-1, -1)$ é um máximo ou mínimo local.
  Como \(f_{xx}(-1, -1) = -6 < 0\)
  o ponto é um ponto de máximo local.
\end{resolution}

%------------------------------------------------------------------------------%

%------------------------------------------------------------------------------%
\addtocounter{exemplo}{1}
\begin{frame}{Exemplo \theexemplo}
  Considerando a função
  \(
    f(x, y) = 4xy -x^4 -y^4
  \)

  a) Calcule o gradiente de $f$

  b) Calcule a hessiana de $f$

  c) Encontre todos os pontos críticos de $f$

  d) Classifique cada ponto crítico de $f$
\end{frame}

%------------------------------------------------------------------------------%
\begin{frame}{Exemplo \theexemplo}
  \noindent\textbf{a)}

  Precisamos das derivadas parciais de $f$
  \[
  \D{f}{x}
  = \D{}{x}\left[4xy -x^4 -y^4\right]
  = 4y - 4x^3
  \]
  \[
  \D{f}{y}
  = \D{}{y}\left[4xy -x^4 -y^4\right]
  = 4x -4y^3
  \]
  então
  \[
  \nabla f = \vetor{4y - 4x^3}{4x -4y^3}
  \]
\end{frame}

%------------------------------------------------------------------------------%
\begin{frame}{Exemplo \theexemplo}
  \noindent\textbf{b)}

  Precisamos das derivadas parciais de segunda ordem de $f$
  \[
  \DS{f}{x}
  = \D{}{x}\left[4y - 4x^3\right]
  = -12x^2
  \]
  \[
  \D{f}{y}
  = \D{}{y}\left[4x -4y^3\right]
  = -12y^2
  \]
  \[
  \DM{f}{x}{y}
  = \D{}{x}\D{f}{y}
  = \D{}{x}\left[4x -4y^3\right]
  = 4
  \]
  então
  \[
  H = \left(\begin{array}{cc}
    -12x^2 & 4 \\
    4 & -12y^2
  \end{array}\right)
  \]
\end{frame}

%------------------------------------------------------------------------------%
\begin{frame}{Exemplo \theexemplo}
  \noindent\textbf{c)}

  A função é um polinômio, então possui derivadas em todos os pontos do plano.
  Assim os pontos críticos são apenas os pontos onde as derivadas parciais são zero,
  \(\nabla f = 0\)
  \[
  4y - 4x^3 = 0
  \qquad\text{e}\qquad
  4x -4y^3 = 0
  \]
  ou, simplificando,
  \[
  y - x^3 = 0
  \qquad\text{e}\qquad
  x - y^3 = 0
  \]
  Isolando $y$ na primeira equação, \(y = x^3\), e substituindo na segunda, temos
  \begin{align*}
    x - y^3                & = 0 \\
    x - \left(x^3\right)^3 & = 0 \\
    x -x^9                 & = 0 \\
    x(1 - x^8)             & = 0
  \end{align*}
  As soluções dessa equação são $x=0$, $x=1$ ou $x-1$.
  Se $x=0$ temos $y=0$,
  se $x=1$ temos $y=1$ e
  se $x=-1$ temos $y=-1$.
  Portanto, os pontos críticos são
  \[
  (x_1, y_1) = (0, 0),
  \qquad\qquad
  (x_2, y_2) = (1, 1)
  \qquad\text{e}\qquad
  (x_3, y_3) = (-1, -1)
  \]
\end{frame}

%------------------------------------------------------------------------------%
\begin{frame}{Exemplo \theexemplo}
  \noindent\textbf{d)}

  Para classificar os pontos críticos precisamos avaliar o discriminante nos
  pontos críticos.
\end{frame}

%------------------------------------------------------------------------------%
\begin{frame}{Exemplo \theexemplo}
  Considerando o ponto \((x_1, y_1) = (0, 0)\)
  \[
  f_{xx}(0, 0) = 0
  \]
  \[
  f_{yy}(0, 0) = 0
  \]
  \[
  f_{xy}(0, 0) = 4
  \]
  \[
  D_1
  = f_{xx}(0, 0)f_{yy}(0, 0) - f^2_{xy}(0, 0)
  = 0 \times 0 - 4^2
  = -16 < 0
  \]
  Portanto, o ponto $(0, 0)$ é um ponto de sela.
\end{frame}

%------------------------------------------------------------------------------%
\begin{frame}{Exemplo \theexemplo}
  Considerando o ponto \((x_2, y_2) = (1, 1)\)
  \[
  f_{xx}(1, 1) = -12
  \]
  \[
  f_{yy}(1, 1) = -12
  \]
  \[
  f_{xy}(1, 1) = 4
  \]
  \[
  D_2
  = f_{xx}(1, 1)f_{yy}(1, 1) - f^2_{xy}(1, 1)
  = (-12)(-12) - 4^2
  = 144 -16
  = 128 > 0
  \]
  Portanto, o ponto $(1, 1)$ é um máximo ou mínimo local. Como \(f_{xx}(1, 1) = -12 < 0\)
  o ponto é um ponto de máximo local.
\end{frame}

%------------------------------------------------------------------------------%
\begin{frame}{Exemplo \theexemplo}
  Considerando o ponto \((x_3, y_3) = (-1, -1)\)
  \[
  f_{xx}(-1, -1) = -12
  \]
  \[
  f_{yy}(-1, -1) = -12
  \]
  \[
  f_{xy}(-1, -1) = 4
  \]
  \[
  D_3
  = f_{xx}(-1, -1)f_{yy}(-1, -1) - f^2_{xy}(-1, -1)
  = (-12)(-12) - 4^2
  = 144 -16
  = 128 > 0
  \]
  Portanto, o ponto $(-1, -1)$ é um máximo ou mínimo local. Como \(f_{xx}(-1, -1) = -12 < 0\)
  o ponto é um ponto de máximo local.
\end{frame}

%------------------------------------------------------------------------------%
\begin{question}
  Considerando a função
  \(
    f(x, y) = x^4 + y^4 + 4xy
  \)

  a) Calcule o gradiente de $f$

  b) Calcule a hessiana de $f$

  c) Encontre todos os pontos críticos de $f$

  d) Classifique cada ponto crítico de $f$
\end{question}

%------------------------------------------------------------------------------%
\begin{resolution}{a}

  Precisamos das derivadas parciais de $f$
  \[
    \D{f}{x}
    = \D{}{x}\left[x^4 + y^4 + 4xy\right]
    = 4x^3 + 4y
  \]
  \[
    \D{f}{y}
    = \D{}{y}\left[x^4 + y^4 + 4xy\right]
    = 4y^3 + 4x
  \]
  então
  \[
    \nabla f = \vetor{4x^3 + 4y}{4y^3 + 4x}
  \]
\end{resolution}

%------------------------------------------------------------------------------%
\begin{resolution}{b}

  Precisamos das derivadas parciais de segunda ordem de $f$
  \[
    \DS{f}{x}
    = \D{}{x}\left[4x^3 + 4y\right]
    = 12x^2
  \]
  \[
    \D{f}{y}
    = \D{}{y}\left[4y^3 + 4x\right]
    = 12y^2
  \]
  \[
    \DM{f}{x}{y}
    = \D{}{x}\D{f}{y}
    = \D{}{x}\left[4y^3 + 4x\right]
    = 4
  \]
  então
  \[
    H = \left(\begin{array}{cc}
      12x^2 & 4 \\
      4 & 12y^2
    \end{array}\right)
  \]
\end{resolution}

%------------------------------------------------------------------------------%
\begin{resolution}{c}
  A função é um polinômio, então possui derivadas em todos os pontos do plano.
  Assim os pontos críticos são apenas os pontos onde as derivadas parciais são zero,
  \(\nabla f = 0\)
  \[
    4x^3 + 4y = 0
    \qquad\text{e}\qquad
    4y^3 + 4x = 0
  \]
  ou, simplificando,
  \[
    x^3 + y = 0
    \qquad\text{e}\qquad
    y^3 + x = 0
  \]
\end{resolution}

%------------------------------------------------------------------------------%
\begin{resolution}{c}
  Isolando $y$ na primeira equação, \(y = -x^3\), e substituindo na segunda, temos
  \begin{align*}
    y^3 + x                 & = 0 \\
    \left(-x^3\right)^3 + x & = 0 \\
    -x^9 + x                & = 0 \\
    x^9 - x                 & = 0 \\
    x(x^8 - 1)              & = 0
  \end{align*}
  As soluções dessa equação são $x=0$, $x=1$ ou $x-1$.
  Se $x=0$ temos $y=0$,
  se $x=1$ temos $y=-1$ e
  se $x=-1$ temos $y=1$.
  Portanto, os pontos críticos são
  \[
    (x_1, y_1) = (0, 0),
    \qquad\qquad
    (x_2, y_2) = (1, -1)
    \qquad\text{e}\qquad
    (x_3, y_3) = (-1, 1)
  \]
\end{resolution}

%------------------------------------------------------------------------------%
\begin{resolution}{d}
  Para classificar os pontos críticos precisamos avaliar o discriminante nos
  pontos críticos.
\end{resolution}

%------------------------------------------------------------------------------%
\begin{resolution}{d}
  Considerando o ponto \((x_1, y_1) = (0, 0)\)
  \[
  f_{xx}(0, 0) = 0
  \]
  \[
  f_{yy}(0, 0) = 0
  \]
  \[
  f_{xy}(0, 0) = 4
  \]
  \[
  D_1
  = f_{xx}(0, 0)f_{yy}(0, 0) - f^2_{xy}(0, 0)
  = 0 \times 0 - 4^2
  = -16 < 0
  \]
  Portanto, o ponto $(0, 0)$ é um ponto de sela.
\end{resolution}

%------------------------------------------------------------------------------%
\begin{resolution}{d}
  Considerando o ponto \((x_2, y_2) = (1, -1)\)
  \[
    f_{xx}(1, -1) = 12
  \]
  \[
    f_{yy}(1, -1) = 12
  \]
  \[
    f_{xy}(1, -1) = 4
  \]
  \[
    D_2
    = f_{xx}(1, -1)f_{yy}(1, -1) - f^2_{xy}(1, -1)
    = 12 \times 12 - 4^2
    = 144 -16
    = 128 > 0
  \]
  Portanto, o ponto $(1, -1)$ é um máximo ou mínimo local.
  Como \(f_{xx}(1, -1) = 12 > 0\)
  o ponto é um ponto de mínimo local.
\end{resolution}

%------------------------------------------------------------------------------%
\begin{resolution}{d}
  Considerando o ponto \((x_3, y_3) = (-1, 1)\)
  \[
    f_{xx}(-1, 1) = 12
  \]
  \[
    f_{yy}(-1, 1) = 12
  \]
  \[
    f_{xy}(-1, 1) = 4
  \]
  \[
    D_3
    = f_{xx}(-1, 1)f_{yy}(-1, 1) - f^2_{xy}(-1, 1)
    = 12 \times 12 - 4^2
    = 144 -16
    = 128 > 0
  \]
  Portanto, o ponto $(-1, 1)$ é um máximo ou mínimo local.
  Como \(f_{xx}(-1, 1) = 12 > 0\)
  o ponto é um ponto de mínimo local.
\end{resolution}

%------------------------------------------------------------------------------%


%------------------------------------------------------------------------------%
\section{Lista mínima}

%------------------------------------------------------------------------------%
\begin{frame}{Lista mínima}
  Cálculo Vol. 2 do Thomas 12\textsuperscript{a} ed. -- Seção 14.7
  \begin{enumerate}
    \item Estudar o texto da seção
    \item Resolver os exercícios: 2, 9, 11, 21, 23, 24, 25, 27
  \end{enumerate}

  \vfill
  Atenção: A prova é baseada no livro, não nas apresentações
\end{frame}

%------------------------------------------------------------------------------%
\end{document}
%------------------------------------------------------------------------------%
